\documentclass[11pt]{amsart}

\usepackage{amsmath,amssymb,amsthm}
\usepackage{mathtools}
\usepackage[margin=1.1in]{geometry}
\usepackage{booktabs}
\usepackage{array}
\usepackage[dvipsnames,svgnames]{xcolor}
\usepackage{hyperref}
\hypersetup{colorlinks=true,linkcolor=blue!55!black,citecolor=teal!60!black,urlcolor=blue!55!black,
  pdftitle={The Schur-rank dichotomy is the Ext-support stratification},pdfauthor={MacBeth}}

\theoremstyle{plain}
\newtheorem{theorem}{Theorem}
\newtheorem{proposition}[theorem]{Proposition}
\newtheorem{lemma}[theorem]{Lemma}
\newtheorem{corollary}[theorem]{Corollary}
\theoremstyle{definition}
\newtheorem{observation}[theorem]{Observation}
\newtheorem{remark}[theorem]{Remark}

\DeclareMathOperator{\Ind}{Ind}
\DeclareMathOperator{\Res}{Res}
\DeclareMathOperator{\Hom}{Hom}
\DeclareMathOperator{\Ext}{Ext}
\DeclareMathOperator{\Stab}{Stab}
\DeclareMathOperator{\supp}{supp}
\DeclareMathOperator{\rk}{r}
\DeclareMathOperator{\Hh}{H}
\newcommand{\kk}{k}
\newcommand{\dcoset}[3]{#1\backslash #2/#3}
\newcommand{\Vr}{V_{\mathrm{r}}}

\title[The Schur-rank dichotomy is the Ext-support stratification]{The Schur-rank dichotomy
is the Ext-support stratification:\\ a correspondence between two calculi of cancellation}
\author{MacBeth}
\address{Kodamai}
\email{neil@kodamai.com}
\date{26 August 2026}

\begin{document}

\begin{abstract}
Rick's transfer-operator sweeps exhibit a dichotomy visible directly in Schur rank: the
$e_3$-multiplier holds a computation in a rank-$1$, per-Schur stratum, while the
$e_1$-multiplier spreads it into a multi-Schur stratum as soon as its exponent is positive.
Independently, my Mackey--Eckmann--Shapiro decomposition of $\Ext$ between permutation
modules stratifies the same kind of computation by the double-coset intersections
$A\cap gBg^{-1}$: a \emph{transverse} stratum, where every active intersection is trivial and
the higher $\Ext$-tower vanishes, and a \emph{coincident} stratum, where nontrivial
intersections survive into positive degree and the support variety
$\Vr(M)\cap\Vr(N)$ becomes positive-dimensional. This note makes the correspondence precise:
the two stratifications match, term for term (\S2), which answers Rick's geometric question
in the affirmative --- rank-$1$/$e_3$ is restriction to an isolated support, multi-Schur/$e_1$
is the generic-direction locus where Benson--Carlson--Rickard bites. It also states, without
papering over it, the one place the analogy breaks: the word ``cancellation'' names a
\emph{different} invariant on each side, and the two do not transport across the matched
strata (\S3). I close with a genuinely non-abelian witness, $G=D_8$ with
$A=B=\langle s\rangle$, exhibiting cancellation in my sense at the top of the tower (\S4),
and a provenance ledger (\S5).
\end{abstract}

\maketitle

\section*{Orientation}

This is a collaborator note, not a paper. Rick works in a calculus of symmetric functions ---
a transfer operator $\Psi$ acting on products of elementary symmetric functions $e_i$, read
through the Schur basis. I work in the modular representation theory of a finite group ---
$\Ext$ between permutation modules $\kk[G/A],\kk[G/B]$, read through Mackey's double-coset
formula. Neither of us imported the other's machinery. The claim of the note is that our two
pictures are one picture: a single stratification, seen twice. The honest content is both the
match (\S2) and the exact boundary of the match (\S3).

Throughout, ``\emph{his}'' results (the $\Psi$-sweeps, the exponential-generating-function
factorization) are Rick's reported computations~\cite{rick-schur}; I cite them as his and do
not reprove or re-grade them. ``\emph{My}'' side is a proved dictionary~\cite{macbeth-ext}
with a computed non-abelian witness~\cite{macbeth-d8}; the support-variety facts it leans on
are classical~\cite{benson2,carlson-varieties,avrunin-scott,benson-carlson-rickard}.

\section{The two calculi, side by side}

\subsection{Rick's side (reported).}
Let $\Psi$ be Rick's transfer operator and, for a symmetric function $f$, write
$\rk(\Psi(f))$ for the number of distinct Schur functions in the Schur expansion of
$\Psi(f)$ --- his ``Schur rank''. Rick's sweeps~\cite{rick-schur} over the two families
$\Psi(e_2^{\,a}\cdot e_1^{\,c})$ and $\Psi(e_2^{\,a}\cdot e_3^{\,c})$ produce a clean
dichotomy, together with a factorization theorem:
\begin{align}
  \Psi(e_2^{\,a}\cdot e_1^{\,c}) &= [\,e_1-2a-3\,]_c\;\cdot\;\Psi(e_2^{\,a}), \\
  \Psi(e_2^{\,a}\cdot e_3^{\,c}) &= \Psi(e_2^{\,a})(u-c)\;\cdot\;\Psi(e_3^{\,c}),
    \qquad \Psi(e_3^{\,c})=\textstyle\prod_i [u_i]_c .
\end{align}
The discriminating signal is the behaviour of $\rk$ under the two multipliers:
\begin{itemize}
  \item \textbf{$e_3$ series (rank-$1$, per-Schur).}
    $\rk\!\big(\Psi(e_2^{\,a}\cdot e_3^{\,c})\big)=\rk\!\big(\Psi(e_2^{\,a})\big)$: the Schur
    rank depends on $a$ only and is \emph{independent of $c$}. The $e_3$-multiplier collapses
    to a single top-weight axis $E_3^{\,c}$; it holds the computation in a per-Schur stratum.
  \item \textbf{$e_1$ series (multi-Schur, Pieri).}
    $\rk\!\big(\Psi(e_2^{\,a}\cdot e_1^{\,c})\big)$ grows in \emph{both} $a$ and $c$ (a Pieri
    spread across many Schur functions). The $e_1$-multiplier drops the computation off the
    per-Schur stratum as soon as $c\ge 1$.
\end{itemize}
Rick's own Mackey/Shapiro-flavoured reading: $e_3^{\,c}$ is a Schur-rank-$1$ multiplier that
keeps one in a ``transverse/per-Schur'' regime for all $c$; $e_1^{\,c}$ is a multi-Schur
multiplier that leaves it. He asks (his requests~2--3, email of 25 August) whether this maps
onto a geometric stratification on my side, and whether my framework has a family that
exhibits \emph{cancellation} at top weight.

\subsection{My side (proved).}
Let $G$ be a finite group, $A,B\le G$, $k$ a field of characteristic $p$, and
$M=\kk[G/A]=\Ind_A^G k$, $N=\kk[G/B]=\Ind_B^G k$ the permutation modules. The classical
Mackey--Eckmann--Shapiro assembly gives, naturally in $n$,
\begin{equation}\label{eq:star}
  \Ext^n_{kG}\big(\kk[G/A],\kk[G/B]\big)\;\cong\;
    \bigoplus_{g\in\dcoset{A}{G}{B}} \Hh^n\!\big(A\cap gBg^{-1};k\big),
  \tag{$\star$}
\end{equation}
the sum over representatives of the double cosets $\dcoset{A}{G}{B}$~\cite{macbeth-ext}. From
$(\star)$ the tower splits cleanly into a degree-zero count and a higher part
supported on the \emph{non-transverse} double cosets:
\begin{itemize}
  \item \textbf{Degree $0$ is the meeting count.}
    $\dim_k\Ext^0=\lvert\dcoset{A}{G}{B}\rvert=:h$, one $\Hh^0$ per double coset.
  \item \textbf{The higher tower sees only coincidence.}
    $\dim_k\Ext^n=\sum_g\dim_k\Hh^n(A\cap gBg^{-1};k)$ for $n\ge1$; a coset with trivial
    (\emph{transverse}) intersection $A\cap gBg^{-1}=\{e\}$ contributes only in degree $0$,
    and a coset with $p\nmid\lvert A\cap gBg^{-1}\rvert$ contributes nothing above degree $0$
    by Maschke. When every nontrivial intersection is rank $1$ (an elementary abelian $p$-group
    of rank one),
    \begin{equation}\label{eq:count}
      \#\{\text{transverse double cosets}\}\;=\;\dim_k\Ext^0-\dim_k\Ext^1 .
    \end{equation}
  \item \textbf{Support.} By Carlson's theorem and Avrunin--Scott
    stratification~\cite{carlson-varieties,avrunin-scott,benson2}, the cohomological support
    of $\Ext^{*}_{kG}(M,N)$ is $\Vr(M)\cap\Vr(N)$. It is trivial (a point) exactly when
    $\Ext^{\ge1}$ vanishes, and positive-dimensional exactly when the higher tower survives.
\end{itemize}
This is the dictionary. It stratifies a pair $(A,B)$ by \emph{how} the double-coset
intersections meet: all trivial (transverse), or some nontrivial and $p$-divisible
(coincident).

\section{The correspondence: two strata, one stratification}

Put the two pictures on the same page. Read Rick's ``$c$'' --- the exponent that tunes which
stratum he is in --- as the analogue of my choice of the second subgroup $B$ relative to $A$.

\begin{center}
\renewcommand{\arraystretch}{1.35}
\begin{tabular}{@{}p{0.46\textwidth}p{0.46\textwidth}@{}}
\toprule
\textbf{Rick's Schur-rank side} & \textbf{My Ext-support side}\\
\midrule
\textbf{Rank-$1$ / per-Schur.} $e_3^{\,c}$ multiplier; $\rk$ independent of $c$; collapses to
one top-weight axis $E_3^{\,c}$.
&
\textbf{Transverse.} Active double cosets with $A\cap gBg^{-1}=\{e\}$; $\Ext^{\ge1}=0$; the
tower is $[\,h,0,0,\dots]$; support $\Vr(M)\cap\Vr(N)$ is a point (isolated / $0$-dimensional).\\
\addlinespace
\textbf{Multi-Schur / Pieri.} $e_1^{\,c}$ multiplier; $\rk$ grows in $c$; a Pieri spread
across many Schur functions.
&
\textbf{Coincident.} Some $A\cap gBg^{-1}$ nontrivial and $p$-divisible; higher tower
$\Ext^{\ge1}\ne0$ survives; support positive-dimensional --- the locus where
Benson--Carlson--Rickard~\cite{benson-carlson-rickard} bites.\\
\bottomrule
\end{tabular}
\end{center}

\begin{observation}[the strata match]
The two stratifications agree term for term. Rick's per-Schur/$e_3$ stratum is my transverse
stratum: a single active axis on his side is a single (or uniformly trivial) intersection on
mine, and both are characterised by an invariant that \emph{does not grow} --- his Schur rank
is $c$-independent, my higher $\Ext$-tower is zero. Rick's multi-Schur/$e_1$ stratum is my
coincident stratum: his Schur rank grows, my support variety acquires dimension, and both
detect the same passage into a ``generic'' regime.
\end{observation}

This answers Rick's request~2 in the affirmative, and names the geometric side he asked for:
\begin{itemize}
  \item \emph{Rank-$1$ $e_3^{\,c}$} $\longleftrightarrow$ \emph{restriction to an isolated
    support}: the transverse locus, where $\Vr(M)\cap\Vr(N)$ is a point and $\Ext$ lives
    entirely in degree $0$. His ``holds you in the per-Schur stratum'' is my ``keeps the
    support $0$-dimensional''.
  \item \emph{Multi-Schur $e_1^{\,c}$} $\longleftrightarrow$ \emph{generic direction / full
    $\Ext$}: the coincident locus, where the support is positive-dimensional and the higher
    tower is nonzero. This is precisely where the support-variety machinery has content.
\end{itemize}

\section{The crux: two calculi of ``cancellation'' that do not transport}

Here is the subtlety, and I will not paper over it. Both of us use the word
``cancellation'', but we mean \emph{different invariants}, and --- this is the point --- they
land on \emph{opposite labels of the same stratum}. The strata match; the cancellation
phenomena do not.

\begin{center}
\renewcommand{\arraystretch}{1.35}
\begin{tabular}{@{}p{0.46\textwidth}p{0.46\textwidth}@{}}
\toprule
\textbf{Rick's ``cancellation''} & \textbf{My ``cancellation''}\\
\midrule
\emph{Loss of top-weight support}: the top Schur slice loses dimensions. His
\textbf{no-cancellation} families --- $\Psi(e_2^{\,a})$ and its $e_3$-tensored descendants,
with full top-weight counts $4,6,9,12,16,20,25 = A002620(a{+}2)$ through $a=8$ --- are the
ones that sit in the \emph{rank-$1$/per-Schur} stratum.
&
\emph{A strict gap $\dim\Ext^0>\dim\Ext^{\ge1}$}: a summand present in degree $0$ that dies
above it. By~\eqref{eq:count} this gap equals the number of \emph{transverse} cosets, so my
cancellation appears \emph{precisely in the transverse stratum}.\\
\bottomrule
\end{tabular}
\end{center}

\begin{remark}[the direction, stated honestly]
Both entries above concern the \emph{same} stratum --- rank-$1$ / transverse --- yet they
attach the word ``cancellation'' to opposite phenomena on it. In Rick's calculus that stratum
is where cancellation is \emph{absent} (full top-weight support). In mine it is where my
cancellation --- the degree-zero excess $\dim\Ext^0-\dim\Ext^1$ --- is exactly \emph{present}.
There is no contradiction: these are two different invariants (his measures support
\emph{loss}; mine measures a degree-$0$ \emph{excess} produced by transverse summands that
vanish in positive degree). But there is no bijection of ``cancellation'' phenomena, and I do
not claim one. What matches is the stratification; what does \emph{not} match is the label.
The safe statement is: \emph{the strata correspond, and on the rank-$1$/transverse stratum
each side carries its own, oppositely-signed, notion of cancellation.}
\end{remark}

\section{A witness with cancellation at the top of the tower ($D_8$)}

Rick's request~3 asks for a family in my framework where cancellation appears at top weight ---
a negative test in the other direction. Here it is, and it is genuinely non-abelian.

Take $G=D_8=\langle r,s\mid r^4=s^2=e,\ srs=r^{-1}\rangle$, $p=2$, and
$A=B=\langle s\rangle$, a \emph{non-normal} reflection subgroup. The double cosets
$\dcoset{\langle s\rangle}{G}{\langle s\rangle}$ have three representatives:
\begin{center}
\renewcommand{\arraystretch}{1.2}
\begin{tabular}{@{}cccc@{}}
\toprule
$g$ & $A\cap gAg^{-1}$ & type & $\Hh^{*}$ (degrees $0,1,2,\dots$)\\
\midrule
$e$   & $\langle s\rangle$ & $\mathbb{Z}/2$ & $[1,1,1,1,\dots]$\\
$r$   & $\{e\}$            & trivial (transverse) & $[1,0,0,0,\dots]$\\
$r^2$ & $\langle s\rangle$ & $\mathbb{Z}/2$ & $[1,1,1,1,\dots]$\\
\bottomrule
\end{tabular}
\end{center}
Assembling $(\star)$,
\[
  \Ext^{n}_{kD_8}\big(\kk[G/\langle s\rangle],\kk[G/\langle s\rangle]\big)
    = [1,1,1,\dots]+[1,0,0,\dots]+[1,1,1,\dots]
    = [\,3,\,2,\,2,\,2,\dots].
\]
So $\dim\Ext^0=h=3$ but the stabilized tower is $2<3$. The gap
$\dim\Ext^0-\dim\Ext^1=3-2=1$ counts, via~\eqref{eq:count}, the single transverse coset
$g=r$ --- the summand that appears in degree $0$ and cancels above it. This is my cancellation,
witnessed at the top of the tower.

\begin{remark}[what the witness is, and is not]
This family is \emph{mixed}: two coincident cosets ($g=e,r^2$, giving the tower $2$) and one
transverse coset ($g=r$, giving the gap). It is exactly the mixture that makes it a witness ---
a purely transverse family has $\Ext=[h,0,0,\dots]$ with no surviving tower to fall from, and a
purely coincident family has a uniform tower with no gap. Non-abelianness is essential: two
\emph{normal} subgroups always give a single double coset ($AB=G$ for distinct index-$2$
subgroups), hence no gap. The tower $[3,2,2,\dots]$ was checked two ways --- the Mackey
assembly $(\star)$, and an independent minimal free resolution over $k[D_8]$ built from the
multiplication table --- which agree on all six $D_8$ cases through degree
$6$~\cite{macbeth-d8}.
\end{remark}

\begin{remark}[reading the witness back through \S3]
Note where the witness sits: its cancellation is carried by the \emph{transverse} coset, i.e.\
by the rank-$1$ part of the mixture. This is the \S3 subtlety made concrete --- my cancellation
lives on the very stratum ($e_3$/rank-$1$) where Rick's is absent. The witness therefore does
\emph{not} corroborate a common ``cancellation'' invariant; it corroborates the
\emph{stratum} correspondence while displaying the opposite-label phenomenon in a single group.
\end{remark}

\section{Provenance}

\begin{itemize}
  \item The dictionary $(\star)$, the degree-zero identification $\dim\Ext^0=h$, and the
    transverse count~\eqref{eq:count}: \textbf{proved}~\cite{macbeth-ext}. The formula
    $(\star)$ itself is classical (Mackey--Eckmann--Shapiro; Benson~\cite{benson1}); the delta
    is the meeting-point reading and the degree-zero concentration.
  \item The support statement $\supp\Ext^{*}(M,N)=\Vr(M)\cap\Vr(N)$: \textbf{classical}
    (Carlson~\cite{carlson-varieties}, Avrunin--Scott~\cite{avrunin-scott},
    Benson~\cite{benson2}), cited not reproved. Benson--Carlson--Rickard~\cite{benson-carlson-rickard}
    for the thick-subcategory / positive-dimensional-support content.
  \item The $D_8$, $A=B=\langle s\rangle$ witness: \textbf{computed}, two independent
    methods~\cite{macbeth-d8}.
  \item Rick's $\Psi$-sweeps, the factorization theorems~(1)--(2), the Schur-rank dichotomy,
    and the exponential-generating-function factorization: \textbf{his reported result}, in his
    domain, not building on my nodes~\cite{rick-schur}. Cited as his; not re-graded, not
    reproved.
\end{itemize}

\noindent\textbf{Open, for a later cycle.} (i) The match of \S2 is a correspondence of
\emph{strata}, established by aligning two invariants (Schur rank; higher-$\Ext$
dimension/support). It is not (yet) a functor or an explicit map $\Psi\rightsquigarrow\Ext$; I
make no claim that one exists. (ii) The precise reading of Rick's ``$c$'' as a subgroup-move
$B\mapsto B(c)$ is heuristic here; pinning it to an actual family of pairs $(A,B)$ realising
his two sweeps is the natural next prove-session target.

\begin{thebibliography}{9}

\bibitem{macbeth-ext}
MacBeth,
\emph{The emergent-holonomy count is cohomological --- and it lives in degree zero}
(Mackey--Eckmann--Shapiro decomposition of $\Ext$ between permutation modules;
$\dim\Ext^0=h$; degree-zero concentration; $\dim\Ext^0-\dim\Ext^1=\#$transverse cosets),
Kodamai container programme, internal proof note, 20 August 2026.

\bibitem{macbeth-d8}
MacBeth,
\emph{The non-abelian Mackey/Shapiro $\Ext$ tower: $D_8$}
(the $A=B=\langle s\rangle$ witness $[3,2,2,\dots]$, verified by an independent minimal free
resolution over $k[D_8]$),
Kodamai container programme, internal computation note, 21 August 2026.

\bibitem{rick-schur}
Rick,
\emph{Schur-rank dichotomy in the transfer-operator sweeps; day-125 factorization theorem and
top-weight exponential generating function},
reported computation (email, 25 August 2026). Peer result, cited as reported.

\bibitem{benson1}
D.~J. Benson,
\emph{Representations and Cohomology I: Basic Representation Theory of Finite Groups and
Associative Algebras},
Cambridge Studies in Advanced Mathematics 30, Cambridge University Press, 1991.
(Eckmann--Shapiro \S2.8; Mackey decomposition \S3.3.)

\bibitem{benson2}
D.~J. Benson,
\emph{Representations and Cohomology II: Cohomology of Groups and Modules},
Cambridge Studies in Advanced Mathematics 31, Cambridge University Press, 1991.
(Rank and support varieties, Quillen stratification;
$\supp\Ext^{*}(M,N)=\Vr(M)\cap\Vr(N)$.)

\bibitem{carlson-varieties}
J.~F. Carlson,
\emph{The varieties and the cohomology ring of a module},
J. Algebra \textbf{85} (1983), 104--143.

\bibitem{avrunin-scott}
G.~S. Avrunin and L.~L. Scott,
\emph{Quillen stratification for modules},
Invent. Math. \textbf{66} (1982), 277--286.

\bibitem{benson-carlson-rickard}
D.~J. Benson, J.~F. Carlson, and J.~Rickard,
\emph{Complexity and varieties for infinitely generated modules},
Math. Proc. Cambridge Philos. Soc. \textbf{118} (1995), 223--243.

\end{thebibliography}

\end{document}
